\documentclass[11pt]{article}

\usepackage[T1]{fontenc}
\usepackage[utf8]{inputenc}
\usepackage{lmodern}
\usepackage[a4paper,margin=1in]{geometry}
\usepackage{amsmath,amssymb}
\usepackage{booktabs}
\usepackage{array}
\usepackage{tabularx}
\usepackage{xcolor}
\usepackage{microtype}
\usepackage{enumitem}
\usepackage{listings}
\usepackage[hidelinks]{hyperref}

\setlength{\parindent}{0pt}
\setlength{\parskip}{0.55em}
\setlength{\emergencystretch}{3em}

\lstset{basicstyle=\ttfamily\small,columns=fullflexible,keepspaces=true,
  frame=single,framerule=0.3pt,xleftmargin=1em,xrightmargin=1em}

\newcommand{\code}[1]{\texttt{\detokenize{#1}}}
\newcommand{\POQueue}{\mathsf{POQueue}}
\newcommand{\Queue}{\mathsf{Queue}}
\newcommand{\AbsQ}{\mathsf{AbsQ}}
\newcommand{\AbsQz}{\mathsf{AbsQ}_0}
\newcommand{\Scan}{\mathsf{Scan}}
\newcommand{\unverified}[1]{\textcolor{gray}{[#1]}}

\title{The Herlihy--Wing Queue Proof in RGSimLin\\
  Compared with Existing Proofs}
\author{Notes for the \code{rbgs-rglin} development}
\date{October 2026}

\begin{document}
\maketitle

\begin{abstract}
We compare the RGSimLin proof of the Herlihy--Wing queue
(\code{examples/HWQ}, written up in \code{HWQueueProofDetails.tex}) with the
major existing proofs of the same algorithm: the mechanized ones in KIV,
Iris/Rocq and TLAPS, and the commonly cited paper proofs.  The closest
relative is the forward-simulation proof of Bouajjani et al.\ (CAV 2017),
whose intermediate specification $\AbsQ$ is essentially our $\POQueue$.  Our
proof differs in where the enqueue's abstract events are placed, which makes
the code-level layer a single-possibility simulation with a plain state
invariant, and in proving the code-independent layer by forward speculation
rather than backward simulation.  Its main weakness is the idealized
underlay.
\end{abstract}

\tableofcontents

\section{Our proof in brief}

The implementation is the standard one over an atomic counter \code{back}
and an unbounded swap array \code{items}:
\begin{lstlisting}
enq(x):  i <- inc(back); store(items[i], x); return ok
deq():   while true:
           r <- read(back);
           for i in 0 .. r-1:
             x <- swap(items[i], null);
             if x != null: return x
\end{lstlisting}
The proof is the vertical composition
\[
  \Queue \;\lhd\; \POQueue \;\lhd\; (\mathsf{Counter} \otimes \mathsf{SwapArray}).
\]

\paragraph{The partial-order queue.}
A $\POQueue$ state is a DAG of enqueued vertices with a set of pending
enqueues and a set of removed vertices.  An enqueue is interval sequential:
its abstract invocation adds a fresh vertex ordered after every
\emph{complete} vertex, and its abstract response marks the vertex complete.
A dequeue is atomic and removes a minimal, live, complete vertex.

\paragraph{Layer 1 (code to $\POQueue$).}
The linearization points are fixed: \code{inc} is the abstract enqueue
invocation, \code{store} is the abstract enqueue response, and a successful
\code{swap} is the dequeue.  The proof is a forward simulation with a single
abstract possibility.  The global invariant says that the counter equals the
next vertex identifier and that slot $i$ holds $x$ exactly when vertex $i$
is live, complete and has value $x$.  The inner loop invariant is
\[
\Scan(\alpha, i, r) \;:\Leftrightarrow\;
\ldots \wedge \forall m < i.\ \mathsf{live}(m) \wedge \mathsf{complete}(m)
  \Rightarrow \forall n.\ (n, m) \in E \Rightarrow n \ge r,
\]
a state predicate that is stable under interference.  It is what shows that
a successful swap takes a minimal vertex.  (\code{HWQueueProof.v}, 1279
lines.)

\paragraph{Layer 2 ($\POQueue$ to $\Queue$).}
The FIFO queue is reached by speculation over sets of possibilities: every
linear extension of the DAG is kept alive, and a dequeue filters the set
down to the extensions headed by the removed vertex.  This layer never
mentions the code.  (\code{QueueLinProof.v}, 1991 lines.)

\paragraph{Result.}
\code{HWQueueTopLevel.MHWQueue_linearizable} is a compositional
linearizability layer theorem against the FIFO queue.  It depends only on
\code{classic} and \code{eq_rect_eq}.  The whole development, including the
specifications, is 4261 lines.

\section{Existing proofs}

Table~\ref{tab:overview} summarizes the proofs we consider.  Sizes are only
indicative: the provers, libraries and proof styles differ, and some
developments were partly written with LLM assistance.

\begin{table}[ht]
\centering\small
\renewcommand{\arraystretch}{1.25}
\begin{tabularx}{\textwidth}{@{}>{\raggedright\arraybackslash}p{3.2cm}
  >{\raggedright\arraybackslash}p{1.9cm}
  >{\raggedright\arraybackslash}X
  >{\raggedright\arraybackslash}p{3.1cm}@{}}
\toprule
Work & Checked in & Technique & Size \\
\midrule
Herlihy, Wing~\cite{HW90} & paper & partial order as auxiliary data,
  Owicki--Gries style, distinct values assumed & --- \\
Schellhorn, Wehrheim, Derrick~\cite{SWD12,SDW14} & KIV & one monolithic
  backward simulation, no auxiliary data & abstract-spec theory alone: 59
  lemmas, 2720 proof steps \\
Henzinger, Sezgin, Vafeiadis~\cite{HSV13,CHSV15} & paper; Cave on a
  linked-list variant & queue aspects (VFresh, VRepet, VOrd, VWit) & --- \\
Bouajjani, Emmi, Enea, Mutluergil~\cite{BEEM17} & paper & forward
  simulation to $\AbsQ$; $\AbsQ \sqsubseteq \AbsQz$ by backward
  simulation & --- \\
Khyzha, Dodds, Gotsman, Parkinson~\cite{KDGP17} & paper &
  rely-guarantee logic building partial-order abstract histories & --- \\
Jung et al.~\cite{Jung20} & Iris (Coq) & logically atomic triple with
  prophecy variables & $\approx$2763 lines \\
Dang et al.\ (Compass)~\cite{Compass22} & iRC11/Iris & relaxed-memory,
  history-based specifications & \unverified{not checked} \\
Jayanti et al.~\cite{JJYH24} & TLAPS & meta-configuration tracking
  (forward, complete) & 2959 proof obligations \\
Zhang, Gregersen, Tassarotti~\cite{ZGT26} & Iris/Rocq & three proofs
  (aspects, Bouajjani's simulation, meta-configurations), each turned into
  a logically atomic triple & 2004 / $\approx$7450 / $\approx$8600
  lines \\
\midrule
This work & Coq (RGSimLin) & singleton forward simulation to $\POQueue$;
  forward speculation to $\Queue$ & 1279 + 1991 lines; 4261 in total \\
\bottomrule
\end{tabularx}
\caption{Proofs of the Herlihy--Wing queue.}
\label{tab:overview}
\end{table}

\subsection{Paper proofs}

\paragraph{Herlihy and Wing~\cite{HW90}.}
The original paper gives a proof sketch.  According to the KIV project page,
it ``uses auxiliary data (a partial order), Owicki-Gries style reasoning and
assumes without loss of generality, that all queue elements are different'',
the last assumption being ``particularly hard to imitate in a formal
proof''.  The 2015 survey of Dongol and Derrick~\cite{DD15} names, besides
this sketch, exactly two formal proofs: the KIV proof and the aspect-oriented
one.

\paragraph{Henzinger, Sezgin and Vafeiadis~\cite{HSV13,CHSV15}.}
Linearizability of a queue is reduced to the absence of four kinds of
violations, which is necessary and sufficient for purely blocking
implementations; the Herlihy--Wing dequeue qualifies because its blocking
executions do not modify the state.  VWit is trivial since a dequeue never
returns empty, and the ordering aspect uses a prophecy instrumentation.  The
proof of the array-based queue is on paper.  The automated check in Cave
covers only a linked-list version, since Cave does not support arrays.

\paragraph{Bouajjani, Emmi, Enea and Mutluergil~\cite{BEEM17}.}
The queue is proved by a forward simulation to an intermediate
specification $\AbsQ$ that keeps a partial order of enqueues.  Dequeues have
fixed linearization points; enqueues have none.  A separate theorem shows
that $\AbsQ$ and the atomic queue $\AbsQz$ refine each other ($\AbsQ
\sqsubseteq \AbsQz$ by a backward simulation, the converse by a forward
one), and a corollary states that a queue with fixed dequeue points refines
$\AbsQz$ exactly when it forward-simulates $\AbsQ$.  The paper assumes that
the same value is never enqueued twice.  Section~\ref{sec:absq} compares
$\AbsQ$ with $\POQueue$ in detail.

\paragraph{Khyzha, Dodds, Gotsman and Parkinson~\cite{KDGP17}.}
A rely-guarantee program logic in which a proof builds an abstract history
whose order is an interval order: an event is added at invocation after all
completed events, edges may be added later, and the history is abstract if
all sequential extensions of its completed part satisfy the specification.
The Herlihy--Wing queue is one of the case studies, in an appendix.
\unverified{We did not retrieve that appendix, so its details are not
compared here.}

\subsection{Mechanized proofs}

\paragraph{KIV~\cite{SWD12,SDW14}.}
The first mechanized proof.  It instantiates a complete refinement theory:
a single backward simulation, built without auxiliary data, relates the
queue to an atomic algebraic queue.  The model is a program with atomic
statements over an infinite array.  The project page lists 59 proved lemmas
with 2720 proof steps and 1265 interactions for the abstract-specification
theory; we found no figure for the whole proof.

\paragraph{Iris, with prophecy variables~\cite{Jung20}.}
The Coq formalization of prophecy variables in Iris includes a proof of
logical atomicity for the queue.  The paper body does not describe this
proof.  The code (a copy of which is in~\cite{ZGT26}) is HeapLang over a
real heap: a \emph{bounded} array fixed at creation, \code{FAA} on the
counter, an enqueue that diverges when the array is full, and a dequeue that
loads a slot and then clears it with a CAS.  Zhang et al.\ give 2763 lines
for this proof.  \unverified{We could not fetch the Iris source itself
(rate-limited), so we do not describe its invariant.}

\paragraph{Compass~\cite{Compass22}.}
A framework extending the iRC11 separation logic for relaxed memory.  Its
abstract reports the first mechanized verification of a relaxed-memory
Herlihy--Wing queue.  This is the only proof under a weak memory model.
\unverified{Details not checked beyond the abstract.}

\paragraph{TLAPS, meta-configuration tracking~\cite{JJYH24}.}
The method keeps the set of all forward simulations to the atomic object
alive and discards one when it fails to extend; it is universal, sound and
complete.  The queue proof is 2959 TLAPS obligations.  The metatheory was
not mechanized in that work; Kent et al.~\cite{KYJBB25} later mechanized
its soundness and completeness in Rocq, with a register as the case study.

\paragraph{Iris/Rocq, three proofs~\cite{ZGT26}.}
A July 2026 preprint proves that logical atomicity is complete for
linearizability, and applies it to the Herlihy--Wing queue three times, on
the Iris code: with the aspect theorem (2004 lines, excluding the
meta-theorem), with Bouajjani et al.'s forward simulation to $\AbsQ$ (4368
lines on the code side, plus 3079 lines for $\AbsQ$ and $\AbsQ
\sqsubseteq \AbsQz$ by backward simulation), and with meta-configuration
tracking (about 8600 lines).  None of these three proofs uses prophecy
variables; the general completeness construction does.  The authors report
that substantial parts of the proofs were automated with LLMs.  The paper
also responds to Hatti et al.~\cite{HOWFS26}, quoting its remark that
``there is no evidence that logical atomicity can be used to specify all
Herlihy-Wing linearizable objects'', and claims to close that gap.

\section{$\POQueue$ versus $\AbsQ$}
\label{sec:absq}

The two intermediate specifications have the same shape.  Both record the
enqueues not yet dequeued as a partial order, add a new enqueue after the
completed ones, and let a dequeue remove a minimal enqueue.  They differ in
where the enqueue's two abstract events happen, and that difference
propagates to the code-level proof.

\begin{table}[ht]
\centering\small
\renewcommand{\arraystretch}{1.25}
\begin{tabularx}{\textwidth}{@{}>{\raggedright\arraybackslash}p{3.4cm}
  >{\raggedright\arraybackslash}X
  >{\raggedright\arraybackslash}X@{}}
\toprule
& $\AbsQ$~\cite{BEEM17} & $\POQueue$ (ours) \\
\midrule
New enqueue vertex & at the method call & at \code{inc} \\
Vertex becomes complete & at the method return & at \code{store} \\
Resulting order & real-time order between enqueues & ``stored before
  incremented'', which contains the real-time order \\
Dequeue & removes a minimal enqueue with the returned value, fixed at the
  swap & removes a minimal, live, complete vertex at the swap \\
Empty dequeue & allowed when every remaining enqueue is pending & no rule;
  a dequeue on an empty queue never responds \\
Identity of an enqueue & operation identifier; distinct values assumed &
  slot index; no assumption on values \\
Removed vertices & dropped from the order & kept as garbage; edges into
  garbage are harmless \\
\bottomrule
\end{tabularx}
\caption{The two intermediate queue specifications.}
\label{tab:absq}
\end{table}

\paragraph{Consequence for the code-level layer.}
Because $\AbsQ$ orders enqueues by real time, the abstract order is not a
function of the concrete state.  Bouajjani et al.'s simulation relation has
to choose an order satisfying three side conditions: pending enqueues are
maximal, the order agrees with the order of reserved slots, and a
condition about dequeues that observed a null slot.  In the Rocq
mechanization of~\cite{ZGT26}, each in-flight dequeue carries a ghost log of
the offsets it has loaded and of the enqueue that had stored each one, used
to discharge minimality at a successful CAS.

Placing the abstract events at \code{inc} and \code{store} makes the order
stronger and determined by the execution: slot $i$ holds a value exactly
when vertex $i$ is live and complete.  The code-level proof is then a forward
simulation with a single possibility, and the observation log is replaced by
$\Scan$, a predicate on the current state that is stable under
interference.  The price is that $\POQueue$ is interval sequential, so the
framework must support linearizability against interval-sequential
specifications and its vertical composition.  This is part of what RGSimLin
provides; $\AbsQ$ is used inside a single simulation and has no such
requirement.

\paragraph{Consequence for the code-independent layer.}
Both proofs separate the future-dependent choice from the code.  Bouajjani
et al.\ prove $\AbsQ \sqsubseteq \AbsQz$ on paper by a backward simulation,
which~\cite{ZGT26} mechanizes in about 3000 lines.  Our Layer~2 proves the
corresponding step forward, over sets of possibilities, in 1991 lines plus
581 lines for $\POQueue$ and its well-formedness lemma.  Neither layer of our
proof uses prophecy variables, backward simulation or history ghost state.
Among the other proofs, only the TLAPS proof~\cite{JJYH24} is forward
throughout, and it is a single proof rather than two layers.

\paragraph{What we do not have.}
Bouajjani et al.'s corollary gives a completeness statement for $\AbsQ$:
every queue with fixed dequeue points that refines $\AbsQz$ forward-simulates
$\AbsQ$.  We have no corresponding statement for $\POQueue$.

\section{Assessment}

\subsection{Distinguishing features of our proof}

\begin{enumerate}[nosep]
\item The future-dependent choice is isolated in a layer that does not
  mention the code, and the code-level layer needs neither possibilities
  nor ghost history.  Bouajjani et al.\ have the same split on paper; our
  placement of the abstract events makes the code-level layer simpler.
\item Reasoning is forward everywhere: no prophecy variables, no backward
  simulation.
\item The result is a compositional linearizability layer theorem that
  composes vertically and horizontally with other layers in the framework.
  The Iris proofs give logically atomic triples, which compose inside Iris;
  the KIV and TLAPS proofs establish linearizability of a closed model.
\item No assumption that enqueued values are distinct.  The original proof
  and Bouajjani et al.\ assume it; the KIV proof also avoids it.
\end{enumerate}

\subsection{Limitations}

\begin{enumerate}[nosep]
\item \textbf{Idealized underlay.}  The bottom layer is the specification of
  an atomic counter and an unbounded atomic swap array.  The Iris proofs are
  about HeapLang code over a real heap, with a bounded array and a
  load-then-CAS dequeue.  Unlike the TS stack development, the queue is not
  yet composed down to concrete memories.
\item \textbf{Sequential consistency.}  Only Compass treats a weak memory
  model.
\item \textbf{No empty result.}  Our dequeue never responds on an empty
  queue.  This matches the algorithm and every proof above; the queue
  specification of~\cite{ZGT26} permits an empty result, but the code never
  produces one.
\item \textbf{No progress.}  Like all the proofs above, ours is a safety
  result.
\end{enumerate}

\subsection{Suggestions}

\begin{itemize}[nosep]
\item Cite $\AbsQ$~\cite{BEEM17} as the precursor of $\POQueue$ in
  \code{HWQueueProofDetails.tex} and in the paper, and state the difference
  in event placement explicitly.
\item The closest quantitative comparison is our Layer~2 against
  $\AbsQ \sqsubseteq \AbsQz$ in~\cite{ZGT26} (forward speculation versus
  backward simulation, about 2600 versus 3000 lines).  A comparison of the
  code-level layers is not meaningful until our queue is composed down to a
  concrete memory.
\end{itemize}

\subsection{Caveats on sources}

The Iris source of~\cite{Jung20} could not be fetched; its code was read
from the copy in the artifact of~\cite{ZGT26}.  The Herlihy--Wing appendix
of~\cite{KDGP17} was not retrieved.  The Compass entry rests on its
abstract.  \cite{ZGT26} is an unrefereed preprint.  The TLAPS proof files
were not read.

\begin{thebibliography}{99}

\bibitem{HW90}
M.~P. Herlihy and J.~M. Wing.
\newblock Linearizability: a correctness condition for concurrent objects.
\newblock \emph{ACM TOPLAS} 12(3):463--492, 1990.

\bibitem{SWD12}
G.~Schellhorn, H.~Wehrheim, and J.~Derrick.
\newblock How to prove algorithms linearisable.
\newblock In \emph{CAV}, 2012.
\newblock KIV project page:
\url{https://swt.informatik.uni-augsburg.de/swt/projects/Herlihy-Wing-queue.html}.

\bibitem{SDW14}
G.~Schellhorn, J.~Derrick, and H.~Wehrheim.
\newblock A sound and complete proof technique for linearizability of
  concurrent data structures.
\newblock \emph{ACM TOCL} 15(4), 2014.

\bibitem{HSV13}
T.~A. Henzinger, A.~Sezgin, and V.~Vafeiadis.
\newblock Aspect-oriented linearizability proofs.
\newblock In \emph{CONCUR}, 2013.

\bibitem{CHSV15}
S.~Chakraborty, T.~A. Henzinger, A.~Sezgin, and V.~Vafeiadis.
\newblock Aspect-oriented linearizability proofs.
\newblock \emph{LMCS} 11(1), 2015.
\newblock \url{https://arxiv.org/abs/1502.07639}.

\bibitem{BEEM17}
A.~Bouajjani, M.~Emmi, C.~Enea, and S.~O. Mutluergil.
\newblock Proving linearizability using forward simulations.
\newblock In \emph{CAV}, 2017.
\newblock \url{https://arxiv.org/abs/1702.02705}.

\bibitem{KDGP17}
A.~Khyzha, M.~Dodds, A.~Gotsman, and M.~Parkinson.
\newblock Proving linearizability using partial orders.
\newblock In \emph{ESOP}, 2017.
\newblock \url{https://arxiv.org/abs/1701.05463}.

\bibitem{Jung20}
R.~Jung, R.~Lepigre, G.~Parthasarathy, M.~Rapoport, A.~Timany, D.~Dreyer,
  and B.~Jacobs.
\newblock The future is ours: prophecy variables in separation logic.
\newblock \emph{PACMPL} 4(POPL), 2020.
\newblock \url{https://plv.mpi-sws.org/prophecies/}.

\bibitem{Compass22}
H.-H. Dang, J.~Jung, J.~Choi, D.-T. Nguyen, W.~Mansky, J.~Kang, and
  D.~Dreyer.
\newblock Compass: strong and compositional library specifications in
  relaxed memory separation logic.
\newblock In \emph{PLDI}, 2022.
\newblock \url{https://plv.mpi-sws.org/compass}.

\bibitem{JJYH24}
P.~Jayanti, S.~Jayanti, U.~Y. Yavuz, and L.~Hern\'andez Videa.
\newblock A universal, sound, and complete forward reasoning technique for
  machine-verified proofs of linearizability.
\newblock \emph{PACMPL} 8(POPL), 2024.
\newblock Artifact:
\url{https://github.com/uguryavuz/machine-certified-linearizability}.

\bibitem{KYJBB25}
Z.~Kent, U.~Y. Yavuz, S.~Jayanti, S.~Balzer, and G.~Blelloch.
\newblock Mechanized metatheory of forward reasoning for end-to-end
  linearizability proofs.
\newblock arXiv:2509.06872, 2025.

\bibitem{ZGT26}
Z.~Zhang, S.~O. Gregersen, and J.~Tassarotti.
\newblock Completeness of logical atomicity for linearizability in
  concurrent separation logic.
\newblock arXiv:2607.11435, 2026.
\newblock Artifact: \url{https://github.com/u8cat/iris-logatom}.

\bibitem{HOWFS26}
E.~Hatti, A.~Oliveira Vale, Z.~Wang, Y.~Feng, and Z.~Shao.
\newblock A complete program logic for compositional linearizability.
\newblock In \emph{ECOOP}, LIPIcs 372, 2026.

\bibitem{DD15}
B.~Dongol and J.~Derrick.
\newblock Verifying linearisability: a comparative survey.
\newblock \emph{ACM Computing Surveys} 48(2), 2015.
\newblock \url{https://arxiv.org/abs/1410.6268}.

\end{thebibliography}

\end{document}
